\section{Empirical Design}

\subsection{Data Sources and Sample Construction}

We combine several datasets to construct the sample. Stock return and trading volume data are obtained from the CRSP Daily Stock File. Firm characteristics are drawn from Compustat. We link CRSP and Compustat using the CRSP--Compustat Merged (CCM) link table. Institutional ownership data are obtained from Thomson Reuters 13F filings via WRDS. Auditor change events are identified from Form 8-K filings (Item 4.01) retrieved from EDGAR. Political data are constructed from publicly available election returns at the county and state levels.

The sample consists of U.S. publicly traded firms with available CRSP and Compustat data. We require each firm-event observation to have valid event dates, return data, and sufficient pre-event observations for estimation. We exclude observations with confounding major announcements (e.g., earnings releases) within the event window where necessary.

\subsection{Identification of Auditor Change Events}

We identify auditor changes using Form 8-K disclosures under Item 4.01 (Changes in Registrant’s Certifying Accountant). For each event, we record the filing date and classify the nature of the change.

We construct the following indicators:

\begin{itemize}
    \item \textbf{Dismissal}: Indicator for auditor dismissal (as opposed to resignation)
    \item \textbf{Big4Switch}: Indicator for switches involving Big 4 auditors
    \item \textbf{Disagreement}: Indicator for disclosures mentioning disagreements or reportable events
\end{itemize}

These classifications allow us to examine heterogeneity in signal ambiguity and credibility.

\subsection{Event Study Methodology}

We compute cumulative abnormal returns (CAR) using a standard market model. Expected returns are estimated over an estimation window of $[-250,-30]$ trading days relative to the event date.

The primary event window is $[-1,+1]$, with robustness checks using alternative windows such as $[-2,+2]$.

\begin{equation}
CAR_{i,t} = \sum_{\tau=-1}^{+1} (R_{i,\tau} - \hat{R}_{i,\tau})
\end{equation}

\subsection{Dependent Variables}

We consider three main outcome variables:

\begin{itemize}
    \item \textbf{CAR}: Directional market reaction
    \item \textbf{Absolute CAR}: $|CAR|$, capturing dispersion in beliefs
    \item \textbf{Abnormal Volume}: 
    \[
    AVOL_{i,t} = \frac{Volume_{i,t}}{\text{Average Volume}_{[-60,-10]}}
    \]
\end{itemize}

Absolute CAR and abnormal volume proxy for disagreement and heterogeneous interpretation.

\subsection{Political Polarization Measures}

\paragraph{Local Polarization}

We measure polarization using election outcomes in the firm's headquarters location. Our primary measure is:

\begin{equation}
Polarization_{local} = |VoteShare_{Dem} - VoteShare_{Rep}|
\end{equation}

Higher values reflect stronger partisan divergence.

\paragraph{Investor-Based Polarization}

Using 13F data, we construct firm-level investor polarization:

\begin{enumerate}
    \item Identify institutional investors holding the firm
    \item Map investors to geographic locations
    \item Assign political leanings based on local election outcomes
    \item Compute dispersion of political alignment across investors
\end{enumerate}

This measure captures heterogeneity in investor beliefs directly.

\subsection{Control Variables}

We include standard controls:

\begin{itemize}
    \item Firm size (log assets)
    \item Market-to-book ratio
    \item Leverage
    \item Return on assets
    \item Stock return volatility
    \item Institutional ownership
\end{itemize}

We also include industry and year fixed effects.

\subsection{Empirical Specifications}

We estimate:

\begin{equation}
Y_{i,t} = \alpha + \beta_1 Event_{i,t} + \beta_2 Polarization_{i,t} + \beta_3 (Event_{i,t} \times Polarization_{i,t}) + \gamma X_{i,t} + \text{FE} + \varepsilon_{i,t}
\end{equation}

where $Y_{i,t}$ is $CAR$, $|CAR|$, or $AVOL$.

The coefficient of interest, $\beta_3$, captures whether polarization amplifies market reactions.

\subsection{Heterogeneity Tests}

We examine:

\begin{itemize}
    \item Big 4 vs. non-Big 4 switches
    \item Dismissal vs. resignation
    \item Presence of disagreement disclosures
    \item High vs. low investor polarization
\end{itemize}

These tests assess whether polarization effects are stronger when signals are more ambiguous or credibility is uncertain.

\subsection{Identification Strategy}

Identification relies on:

\begin{itemize}
    \item Sharp event timing (8-K disclosures)
    \item Cross-sectional variation in polarization
    \item Fixed effects controlling for unobserved heterogeneity
\end{itemize}

Standard errors are clustered at the firm level.

\subsection{Robustness Tests}

We perform:

\begin{itemize}
    \item Alternative event windows
    \item Alternative polarization measures
    \item Exclusion of overlapping events
    \item Placebo tests using non-audit 8-K filings
\end{itemize}

\subsection{Table Structure}

\begin{itemize}
    \item \textbf{Table 1}: Summary statistics
    \item \textbf{Table 2}: Baseline results
    \item \textbf{Table 3}: Heterogeneity
    \item \textbf{Table 4}: Robustness
\end{itemize}